\chapter{Deret Laurent dan Teorema Residu}
\label{ch:b3:residues}

Apa yang terjadi pada sebuah \omterm{def:b3:holomorphic:holo}{fungsi holomorfik} di dekat titik yang ia
\emph{tak} terdefinisi di sana? Jawabannya berupa trikotomi \omterm{def:b3:complete:complete}{lengkap} ---
titik terhapuskan, kutub, atau kesingularan hakiki --- yang terbaca dari
deret pangkat dua sisi, yakni uraian Laurent. Satu koefisien
uraian itu, yakni \emph{\omterm{def:b3:residues:singularities}{residunya}}, mengendalikan setiap \omterm{def:b3:holomorphic:contour}{integral kontur}
di sekitar kesingularannya: sehingga teorema \omterm{def:b3:residues:singularities}{residu} mengubah
integral tentu yang sukar menjadi aljabar berhingga, mencacah akar
fungsi (lewat asas argumen dan Rouché), serta membuktikan teorema pemetaan
terbuka. Mula-mula kita meningkatkan teorema Cauchy dari
daerah berbentuk bintang ke bentuk definitifnya yang bebas homologi ---
lewat argumen elegan Dixon --- sehingga semua kontur yang berbilangan
lilitan nol di sekitar komplemennya menjadi tersedia.

\section{Teorema Cauchy global}

Sebuah \emph{siklus}\index{siklus} $\Gamma$ adalah jumlah formal berhingga
atas \omterm{def:b3:topology:pathconnected}{lintasan} tertutup $\gamma_1, \dots, \gamma_m$; sedangkan integral dan
indeks sepanjang $\Gamma$ adalah jumlah yang bersesuaian, dan
$\operatorname{im}\Gamma = \bigcup\operatorname{im}\gamma_j$.

\begin{theorem}[Cauchy, bentuk global]
\label{thm:b3:residues:globalcauchy}
Misalkan $\Omega \subseteq \C$ terbuka, $f \in \mathcal H(\Omega)$,
dan $\Gamma$ sebuah siklus di $\Omega$ sedemikian sehingga
\[
\operatorname{Ind}_\Gamma(w) = 0
\qquad\text{untuk setiap } w \notin \Omega .
\]
Maka, untuk setiap $z \in \Omega\setminus\operatorname{im}\Gamma$,
\[
\frac1{2\iu\pi}\int_\Gamma\frac{f(w)}{w - z}\,\dd w =
\operatorname{Ind}_\Gamma(z)\,f(z),
\qquad\text{dan}\qquad
\int_\Gamma f(w)\,\dd w = 0 .
\]
\index{teorema Cauchy (global)}
\end{theorem}

\begin{proof}[Bukti (Dixon)]
Definisikan $g \colon \Omega\times\Omega \to \C$ lewat
\[
g(z, w) = \begin{cases}
\dfrac{f(w) - f(z)}{w - z} & w \neq z,\\[4pt]
f'(z) & w = z .
\end{cases}
\]
\emph{Fungsi $g$ \omterm{def:b3:topology:continuity}{kontinu}}: di luar diagonalnya, jelas. Sedangkan di dekat sebuah
titik diagonal $(a, a)$, uraikan $f$ menjadi deret pangkat di $a$
(\cref{thm:b3:holomorphic:analytic}): $f(w) - f(z) =
\sum_{n\geq1}c_n\bigl((w-a)^n - (z-a)^n\bigr)$, lalu membagi
setiap sukunya dengan $w - z$ (lewat pemfaktoran $u^n - v^n$) memberikan, untuk
$z, w \in D(a, r)$,
\[
g(z, w) =
\sum_{n\geq1}c_n\sum_{j=0}^{n-1}(w-a)^{\,j}(z-a)^{\,n-1-j},
\]
yang sahih pula pada diagonalnya (sebab setiap jumlah dalamnya menjadi
$n(z-a)^{n-1}$, yang berjumlah $f'(z)$). Untuk $r$ yang kecil deretnya
konvergen seragam pada $D(a,r)^2$
(sebab $\abs{\text{sukunya}} \leq n\abs{c_n}r^{n-1}$, yang terjumlahkan di dalam
jari-jarinya): sehingga jumlahnya \omterm{def:b3:topology:continuity}{kontinu}.

Tetapkan $h(z) = \frac1{2\iu\pi}\int_\Gamma g(z, w)\,\dd w$ pada
$\Omega$: ia \omterm{def:b3:topology:continuity}{kontinu} (lewat \omterm{def:b3:topology:continuity}{kekontinuan} seragam $g$ pada \omterm{def:b3:topology:compact}{kompak}),
dan \omterm{def:b3:holomorphic:holo}{holomorfik} --- menurut Morera
(yakni kriteria \cref{thm:b3:holomorphic:weierstrassconv}): sebab untuk sebuah
segitiga $T \subseteq \Omega$, Fubini memberikan $\int_{\partial
T}h = \frac1{2\iu\pi}\int_\Gamma\bigl(\int_{\partial T}g(z,
w)\dd z\bigr)\dd w = 0$, dengan integral dalamnya lenyap karena
$z \mapsto g(z, w)$ \omterm{def:b3:holomorphic:holo}{holomorfik} pada $\Omega$ (sebab di $z = w$
kesingularannya terhapuskan: karena $g$ \omterm{def:b3:topology:continuity}{kontinu} di sana dan
\omterm{def:b3:holomorphic:holo}{holomorfik} di tempat lain --- yakni argumen perluasan pada
bukti \cref{thm:b3:holomorphic:formula}).

Pada \omterm{def:b3:topology:topology}{himpunan terbuka} $\Omega' = \{z \notin \operatorname{im}\Gamma
: \operatorname{Ind}_\Gamma(z) = 0\}$, definisikan $h_1(z) =
\frac1{2\iu\pi}\int_\Gamma\frac{f(w)}{w - z}\dd w$: yang
\omterm{def:b3:holomorphic:holo}{holomorfik} pada $\Omega'$ (lewat Morera atau pendiferensialan di bawah
integralnya). Untuk $z \in \Omega\cap\Omega'$:
\[
h(z) = \frac1{2\iu\pi}\int_\Gamma\frac{f(w)}{w-z}\dd w
- f(z)\operatorname{Ind}_\Gamma(z) = h_1(z) .
\]
Menurut hipotesisnya $\Omega\cup\Omega' = \C$ (sebab $w \notin \Omega
\Rightarrow \operatorname{Ind}_\Gamma(w) = 0$), sehingga $h$ dan
$h_1$ merekat menjadi fungsi utuh $H$. Lalu karena komponen tak terbatas
komplemen $\operatorname{im}\Gamma$ terletak
di $\Omega'$ dan $h_1(z) \to 0$ saat $\abs z \to \infty$ (lewat batas
ML), maka $H$ terbatas dan menuju $0$: sehingga Liouville
(\cref{cor:b3:holomorphic:liouville}) memberikan $H \equiv 0$. Jadi
$h \equiv 0$ pada $\Omega$, yang merupakan rumus integralnya.
Lalu dengan menerapkannya, untuk $a \in
\Omega\setminus\operatorname{im}\Gamma$ yang tetap, pada $\tilde f(w) =
(w - a)f(w)$ di $z = a$:
\[
\frac1{2\iu\pi}\int_\Gamma f(w)\dd w =
\frac1{2\iu\pi}\int_\Gamma \frac{\tilde f(w)}{w - a}\dd w =
\operatorname{Ind}_\Gamma(a)\,\tilde f(a) = 0 .
\]
\end{proof}

\section{Deret Laurent dan kesingularan terpencil}

\begin{theorem}[Uraian Laurent]\label{thm:b3:residues:laurent}
Misalkan $f$ \omterm{def:b3:holomorphic:holo}{holomorfik} pada anulus $A = \{r < \abs{z - a} <
R\}$ (dengan $0 \leq r < R \leq \infty$). Maka\index{deret Laurent}
\[
f(z) = \sum_{n\in\Z}c_n\,(z - a)^n
\qquad\text{pada } A,
\qquad
c_n = \frac1{2\iu\pi}\int_{C_\rho}\frac{f(w)}{(w -
a)^{n+1}}\,\dd w
\]
untuk sembarang $r < \rho < R$ (dan tak bergantung pada $\rho$), dengan kedua
separuh deretnya konvergen normal pada subanulus \omterm{def:b3:topology:compact}{kompak}. Sedangkan
uraiannya tunggal.
\end{theorem}

\begin{proof}
Tetapkan $r < \rho_1 < \abs{z - a} < \rho_2 < R$ lalu misalkan $\Gamma =
C_{\rho_2} - C_{\rho_1}$ (dengan yang luar berlawanan arah jarum jam dan yang dalam
searah): yakni sebuah siklus di $A$ dengan
$\operatorname{Ind}_\Gamma(w) = 0$ untuk setiap $w \notin A$
(sebab titik di dalam cakram kecilnya: $1 - 1 = 0$; sedangkan di luar cakram
besarnya: $0 - 0$). Menurut \cref{thm:b3:residues:globalcauchy},
$\operatorname{Ind}_\Gamma(z) = 1 - 0 = 1$ memberikan
\[
f(z) = \frac1{2\iu\pi}\int_{C_{\rho_2}}\frac{f(w)}{w - z}\dd
w - \frac1{2\iu\pi}\int_{C_{\rho_1}}\frac{f(w)}{w - z}\dd w .
\]
Uraikan kernel pertamanya seperti pada
\cref{thm:b3:holomorphic:analytic} (yakni pangkat $\frac{z -
a}{w - a}$ yang bermodulus $< 1$): diperoleh bagian tak negatifnya
$\sum_{n\geq0}c_n(z-a)^n$. Sedangkan pada yang kedua, uraikanlah ke arah
sebaliknya: $\frac{-1}{w - z} = \frac1{(z-a)(1 - \frac{w - a}{z -
a})} = \sum_{m\geq0}\frac{(w-a)^m}{(z - a)^{m+1}}$, yang konvergen
normal pada $C_{\rho_1}$: diperoleh bagian negatifnya
$\sum_{n\leq-1}c_n(z-a)^n$ dengan koefisien yang dinyatakan itu (dengan indeks
$n = -m-1$). Untuk ketakbergantungannya pada $\rho$: integral koefisiennya
atas $C_{\rho}$ dan $C_{\rho'}$ berselisih sebesar
$\int_\Gamma$ atas sebuah siklus berindeks nol bagi
$\frac{f(w)}{(w-a)^{n+1}}$ yang \omterm{def:b3:holomorphic:holo}{holomorfik} di $A$: jadi nol, menurut
\cref{thm:b3:residues:globalcauchy} lagi. Untuk ketunggalannya:
integralkan $\sum c_n(z-a)^n$ terhadap $(z - a)^{-m-1}$ atas
$C_\rho$ suku demi suku (lewat kekonvergenan normalnya): hanya $n = m$
yang bertahan.
\end{proof}

\begin{definition}\label{def:b3:residues:singularities}
Jika $f$ \omterm{def:b3:holomorphic:holo}{holomorfik} pada cakram tertusuk $D(a, R)\setminus
\{a\}$, uraikanlah lewat Laurent (dengan $r = 0$). Ada tiga
kasus yang saling eksklusif:\index{kesingularan terpencil}
\begin{itemize}
  \item semua $c_n = 0$ untuk $n < 0$: yakni \emph{\omterm{thm:b3:residues:riemanncw}{kesingularan
        terhapuskan}} (sebab deret tak negatifnya memperluas $f$
        secara \omterm{def:b3:holomorphic:holo}{holomorfik} ke $a$);
  \item $c_n \neq 0$ untuk berhingga banyak, setidaknya satu, $n < 0$:
        yakni sebuah \emph{kutub} berorde $m = -\min\{n : c_n \neq 0\}$;
        setara dengan itu $f = g/(z-a)^m$ dengan $g$ \omterm{def:b3:holomorphic:holo}{holomorfik} dan $g(a)
        \neq 0$; setara pula dengan $\abs{f(z)} \to \infty$ saat $z\to
        a$;
  \item tak berhingga banyak $c_n \neq 0$ yang negatif: yakni
        \emph{kesingularan hakiki}.
\end{itemize}
Sedangkan \emph{residu}\index{residu} di sana adalah
$\operatorname{Res}(f, a) = c_{-1}$. Dan fungsi yang \omterm{def:b3:holomorphic:holo}{holomorfik} pada
$\Omega$ dikurangi sebuah himpunan kutub disebut
\emph{meromorfik}\index{fungsi meromorfik} pada $\Omega$.
\end{definition}

\begin{theorem}[Riemann; Casorati--Weierstrass]
\label{thm:b3:residues:riemanncw}
Misalkan $f$ \omterm{def:b3:holomorphic:holo}{holomorfik} pada $D(a,R)\setminus\{a\}$.
\begin{enumerate}
  \item (Riemann) Jika $f$ \emph{terbatas} di dekat $a$, maka
        kesingularannya terhapuskan.\index{kesingularan terhapuskan}
  \item (Casorati--Weierstrass) Jika $a$ bersifat hakiki, maka
        $f\bigl(D(a,\varepsilon)\setminus\{a\}\bigr)$ bersifat padat
        di $\C$ untuk setiap $\varepsilon$.
        \index{teorema Casorati--Weierstrass}
\end{enumerate}
\end{theorem}

\begin{proof}
(1) Untuk $n < 0$ dan $\rho \to 0$: $\abs{c_n} \leq
M\rho^{-n-1}\cdot\rho\cdot\rho^{-\,n}\dots$ lewat ML pada
$C_\rho$: sehingga $\abs{c_n} \leq \frac{1}{2\pi}\,2\pi\rho\cdot
M\rho^{-(n+1)} = M\rho^{-n} \to 0$ (sebab $-n > 0$): jadi semua
koefisien negatifnya lenyap.
(2) Seandainya suatu nilai $b$ tak didekati: maka $\abs{f - b} \geq
\delta$ di dekat $a$, sehingga $g = 1/(f - b)$ \omterm{def:b3:holomorphic:holo}{holomorfik} dan
terbatas di dekat $a$: jadi terhapuskan menurut (1), dan $g$ diperluas bernilai $c$.
Jika $c \neq 0$, maka $f = b + 1/g$ terbatas di dekat $a$: jadi terhapuskan
--- yang terkecualikan. Sedangkan jika $c = 0$, maka $g$ berakar berorde berhingga $m$
di $a$ (\cref{thm:b3:holomorphic:identity}; sebab $g \not\equiv 0$),
dan $f = b + 1/g$ berkutub orde $m$: yang terkecualikan lagi.
\end{proof}

\section{Teorema residu}

\begin{theorem}[Teorema residu]\label{thm:b3:residues:residue}
Misalkan $\Omega$ terbuka, $S \subseteq \Omega$ berhingga, $f \in
\mathcal H(\Omega\setminus S)$, dan $\Gamma$ sebuah siklus di
$\Omega\setminus S$ dengan $\operatorname{Ind}_\Gamma(w) = 0$
untuk setiap $w \notin \Omega$. Maka\index{teorema residu}
\[
\frac{1}{2\iu\pi}\int_\Gamma f(z)\,\dd z
= \sum_{a\in S}\operatorname{Ind}_\Gamma(a)\,
\operatorname{Res}(f, a) .
\]
\end{theorem}

\begin{proof}
Untuk setiap $a \in S$, misalkan $P_a(z) = \sum_{n\leq-1}c_n^{(a)}(z -
a)^n$ bagian pokok $f$ di $a$: yakni deret yang
konvergen pada $\C\setminus\{a\}$ (sebab jari-jarinya dalam $1/(z-a)$ bersifat
tak berhingga: karena ekor Laurentnya konvergen untuk setiap $\abs{z-a}$ yang kecil
sehingga, sebagai deret pangkat dalam $(z-a)^{-1}$, di mana-mana), dan
\omterm{def:b3:holomorphic:holo}{holomorfik} di sana. Maka $g = f - \sum_{a\in S}P_a$ berkesingularan
terhapuskan di setiap titik $S$ (sebab uraian Laurentnya
di $a$ tak berbagian negatif: karena $P_{a'}$ yang lain
\omterm{def:b3:holomorphic:holo}{holomorfik} di $a$), sehingga $g$ diperluas secara \omterm{def:b3:holomorphic:holo}{holomorfik} ke
$\Omega$, dan \cref{thm:b3:residues:globalcauchy} memberikan
$\int_\Gamma g = 0$. Lalu tinggal mengintegralkan setiap $P_a$:
suku demi suku (lewat kekonvergenan normal pada
$\operatorname{im}\Gamma$ yang \omterm{def:b3:topology:compact}{kompak}, yang menghindari $a$),
\[
\frac1{2\iu\pi}\int_\Gamma(z - a)^n\,\dd z = 0 \ (n \leq -2:
\text{yakni antiturunan } \tfrac{(z-a)^{n+1}}{n+1}),
\qquad
\frac1{2\iu\pi}\int_\Gamma\frac{\dd z}{z - a} =
\operatorname{Ind}_\Gamma(a),
\]
sehingga $\frac1{2\iu\pi}\int_\Gamma P_a =
c_{-1}^{(a)}\operatorname{Ind}_\Gamma(a)$. Lalu jumlahkan atas $a$.
\end{proof}

\begin{method}[Menghitung residu]\label{met:b3:residues:compute}
Untuk kutub sederhana: $\operatorname{Res}(f, a) = \lim_{z\to a}(z -
a)f(z)$; dan untuk $f = g/h$ dengan $g(a) \neq 0$, $h(a) = 0$, dan $h'(a)
\neq 0$: $\operatorname{Res} = g(a)/h'(a)$. Untuk kutub berorde $m$:
$\operatorname{Res}(f, a) = \frac1{(m-1)!}\lim_{z\to
a}\bigl((z-a)^mf(z)\bigr)^{(m-1)}$. Untuk kesingularan hakikinya:
uraikanlah lalu bacalah $c_{-1}$ (misalnya dari deret yang dikenal). Dan periksalah selalu
kutub mana yang sungguh dilingkari konturnya, dan dengan
indeks yang mana.
\end{method}

\begin{example}[Empat tipe integral klasiknya]
\label{ex:b3:residues:integrals}
(a) \emph{Rasional atas $\R$}: untuk $\int_\R\frac{\dd x}{1 +
x^4}$, tutuplah dengan setengah lingkaran besar $S_R$ di setengah bidang
atasnya: sebab integrannya $O(R^{-4})$ di sana, sehingga
$\int_{S_R} \to 0$ (lewat ML), dan teorema \omterm{def:b3:residues:singularities}{residu} dengan
kutub $\eu^{\iu\pi/4}, \eu^{3\iu\pi/4}$ (yang sederhana, berresidu
$\frac1{4z^3} = \frac{z}{4z^4} = -\frac z4$ di sebuah kutub) memberikan
\[
\int_\R\frac{\dd x}{1 + x^4}
= 2\iu\pi\Bigl(-\frac{\eu^{\iu\pi/4}}4 -
\frac{\eu^{3\iu\pi/4}}4\Bigr)
= \frac{\pi}{\sqrt2} .
\]
(b) \emph{Tipe Fourier}: untuk $t \geq 0$,
$\int_\R\frac{\eu^{\iu tx}}{1 + x^2}\dd x =
2\iu\pi\operatorname{Res}\bigl(\tfrac{\eu^{\iu tz}}{1+z^2},
\iu\bigr) = 2\iu\pi\frac{\eu^{-t}}{2\iu} = \pi\eu^{-t}$ ---
sebab setengah lingkaran atasnya berhasil karena $\abs{\eu^{\iu tz}} =
\eu^{-t\operatorname{Im}z} \leq 1$ di sana; lalu dengan mengambil bagian realnya:
$\int_\R\frac{\cos(tx)}{1+x^2}\dd x = \pi\eu^{-\abs t}$,
yang menuntaskan rumus yang diterima pada \cref{exo:b3:lebesgue:10}.
(c) \emph{Trigonometri atas satu periode}: substitusikan $z =
\eu^{\iu t}$, $\cos t = \frac{z + z^{-1}}2$, dan $\dd t =
\frac{\dd z}{\iu z}$: maka $\int_0^{2\pi}\frac{\dd t}{a + \cos t}$
(dengan $a > 1$) menjadi pencacahan \omterm{def:b3:residues:singularities}{residu} di dalam lingkaran satuannya
(\cref{exo:b3:residues:2}).
(d) \emph{Deret}: pasangkan $f$ dengan $\pi\cot(\pi z)$, yang kutubnya
berupa bilangan bulat berresidu $1$: sehingga soal akhir pekannya menjumlahkan
$\sum n^{-2}$ dan $\sum n^{-4}$ dengan cara ini.
\end{example}

\begin{omfigure}
\begin{tikzpicture}[scale=1.15]
  \draw[->] (-2.5,0) -- (2.6,0) node[right, font=\small]
    {$\operatorname{Re}$};
  \draw[->] (0,-0.5) -- (0,2.5) node[above, font=\small]
    {$\operatorname{Im}$};
  \draw[omThm, very thick, ->] (-2,0) -- (0.3,0);
  \draw[omThm, very thick] (0.3,0) -- (2,0);
  \draw[omThm, very thick] (2,0) arc (0:90:2);
  \draw[omThm, very thick, ->] (2,0) arc (0:50:2);
  \draw[omThm, very thick] (0,2) arc (90:180:2);
  \node[omThm, font=\small] at (1.75,1.55) {$S_R$};
  \node[font=\small] at (2.2,-0.25) {$R$};
  \node[font=\small] at (-2.15,-0.25) {$-R$};
  \fill[omDef] (45:1.19) circle (1.6pt)
    node[above right, font=\small] {$\eu^{\iu\pi/4}$};
  \fill[omDef] (135:1.19) circle (1.6pt)
    node[above left, font=\small] {$\eu^{3\iu\pi/4}$};
  \fill[black!40] (-45:1.19) circle (1.6pt);
  \fill[black!40] (-135:1.19) circle (1.6pt);
\end{tikzpicture}

{\small Kontur setengah lingkaran bagi $\int_\R\frac{\dd x}{1 +
x^4}$: saat $R \to \infty$ busurnya menyumbang $O(R^{-3})$, dan
teorema \omterm{def:b3:residues:singularities}{residunya} mencacah kedua kutub yang terlingkupi (biru). Sedangkan
kedua kutub bawahnya (kelabu) berada di luar: jadi berindeks $0$.}
\end{omfigure}

\section{Asas argumen dan teorema Rouché}

\begin{theorem}[Asas argumen]
\label{thm:b3:residues:argument}
Misalkan $f$ \omterm{def:b3:residues:singularities}{meromorfik} pada $\Omega$, berakar $z_j$ (berorde
$m_j$) dan berkutub $p_k$ (berorde $\mu_k$), dan $\gamma$ sebuah \omterm{def:b3:topology:pathconnected}{lintasan}
tertutup di $\Omega$ yang menghindari semuanya, dengan
$\operatorname{Ind}_\gamma = 0$ di luar $\Omega$. Maka
\index{asas argumen}
\[
\frac1{2\iu\pi}\int_\gamma\frac{f'(z)}{f(z)}\,\dd z
= \sum_j m_j\operatorname{Ind}_\gamma(z_j)
- \sum_k \mu_k\operatorname{Ind}_\gamma(p_k)
\]
(dengan berhingga banyak sukunya yang tak nol). Untuk kontur sederhana
yang berlawanan arah jarum jam, integralnya mencacah akar dikurangi
kutub di dalamnya, beserta kelipatannya --- dan sama dengan bilangan
lilitan \omterm{def:b3:topology:pathconnected}{lintasan} petanya $f\circ\gamma$ di sekitar $0$.
\end{theorem}

\begin{proof}
Di dekat sebuah akar berorde $m$: $f = (z-a)^mg$ dengan $g(a) \neq 0$, sehingga
$\frac{f'}f = \frac m{z - a} + \frac{g'}g$ dengan suku keduanya
\omterm{def:b3:holomorphic:holo}{holomorfik} di dekat $a$: yakni kutub sederhana berresidu $m$. Sedangkan di dekat
kutub berorde $\mu$: $f = (z-a)^{-\mu}g$ memberikan \omterm{def:b3:residues:singularities}{residu}
$-\mu$. Di tempat lain $\frac{f'}f$ bersifat \omterm{def:b3:holomorphic:holo}{holomorfik}. (Sedangkan akar dan
kutub yang berindeks tak nol terletak di daerah \omterm{def:b3:topology:compact}{kompak} yang dilingkari
$\gamma$; dan menurut teorema keidentikannya keduanya berhingga banyak
di sana, sebab $f \not\equiv 0$.) Lalu terapkan
\cref{thm:b3:residues:residue}. Untuk catatan terakhirnya:
$\frac1{2\iu\pi}\int_\gamma\frac{f'}f =
\frac1{2\iu\pi}\int_{f\circ\gamma}\frac{\dd w}w =
\operatorname{Ind}_{f\circ\gamma}(0)$ (substitusikan $w =
f(\gamma(t))$).
\end{proof}

\begin{theorem}[Rouché]\label{thm:b3:residues:rouche}
Misalkan $f, g$ \omterm{def:b3:holomorphic:holo}{holomorfik} pada $\Omega$, dan $\gamma$ sebuah \omterm{def:b3:topology:pathconnected}{lintasan}
tertutup dengan $\operatorname{Ind}_\gamma \in \{0,1\}$ yang nol di luar
$\Omega$ (yakni kontur sederhana). Jika\index{teorema Rouché}
\[
\abs{g(z)} < \abs{f(z)} \qquad \text{pada }
\operatorname{im}\gamma,
\]
maka $f$ dan $f + g$ berakar sama banyak (beserta
kelipatannya) di daerah $\{\operatorname{Ind}_\gamma =
1\}$.
\end{theorem}

\begin{proof}
Untuk $t \in \intcc01$, $f_t = f + tg$ tak berakar pada
$\operatorname{im}\gamma$ (sebab $\abs{f_t} \geq \abs f - \abs g >
0$), sehingga
\[
N(t) = \frac1{2\iu\pi}\int_\gamma
\frac{f_t'(z)}{f_t(z)}\,\dd z
\]
terdefinisi baik; dan ia mencacah akar di daerah yang terlingkupi
(\cref{thm:b3:residues:argument}; sebab tak ada kutub). Lalu $N$ bersifat
\omterm{def:b3:topology:continuity}{kontinu} terhadap $t$ (sebab integrannya \omterm{def:b3:topology:continuity}{kontinu} bersama, dengan
penyebutnya terbatas dari bawah secara seragam --- lewat kekonvergenan
terdominasi) dan bernilai bulat: jadi konstan. Sehingga $N(0) = N(1)$.
\end{proof}

\begin{corollary}[Teorema pemetaan terbuka]
\label{cor:b3:residues:openmapping}
Sebuah \omterm{def:b3:holomorphic:holo}{fungsi holomorfik} tak konstan pada \omterm{def:b3:topology:topology}{himpunan terbuka} \omterm{def:b3:topology:connected}{terhubung} merupakan
pemetaan terbuka. Khususnya (lagi) asas maksimumnya
berlaku, dan bijeksi \omterm{def:b3:holomorphic:holo}{holomorfik} berbalikan
\omterm{def:b3:holomorphic:holo}{holomorfik}.\index{teorema pemetaan terbuka (holomorfik)}
\end{corollary}

\begin{proof}
Misalkan $f(a) = b$; maka $f - b$ berakar berorde berhingga $m
\geq 1$ di $a$ (menurut teorema keidentikannya: sebab $f \not\equiv b$). Pilihlah
$r$ yang $f - b$-nya bebas akar pada $\bar D(a, r)\setminus\{a\}$
(menurut \omterm{thm:b3:holomorphic:identity}{akar terpencilnya}) lalu misalkan $\delta = \min_{\abs{z - a} =
r}\abs{f(z) - b} > 0$. Untuk $\abs{w - b} < \delta$: pada
lingkarannya, $\abs{(b - w)} < \delta \leq \abs{f - b}$, sehingga
Rouché ($f - b$ terhadap konstanta $b - w$) mengatakan bahwa $f - w$ berakar
tepat $m$ kali di $D(a, r)$: sehingga setiap $w$ semacam itu tercapai
--- jadi $f(D(a,r)) \supseteq D(b, \delta)$: yakni terbuka. Untuk asas
maksimumnya: maksimum $\abs f$ di \omterm{def:b3:topology:interior}{interiornya} mustahil bagi
$f$ yang tak konstan, sebab petanya di sekitar $f(a)$ memuat titik
bermodulus lebih besar. Untuk balikannya: bijeksi \omterm{def:b3:holomorphic:holo}{holomorfik} $f$ bersifat
terbuka, sehingga $f^{-1}$ \omterm{def:b3:topology:continuity}{kontinu}; dan akar $f - f(a)$ di
$a$ bersifat sederhana (sebab $m \geq 2$ akan memberikan $m$ prapeta nilai
di dekatnya --- yang berbeda-beda, karena $f'$ lenyap hanya di
titik terpencil, sehingga di dekat $a$ ke-$m$ akar $f - w$ bersifat
sederhana dan berbeda untuk $w$ kecil yang generik: yang bertentangan
dengan keinjektifannya); lalu $f'(a) \neq 0$ dan hasil bagi
selisih $f^{-1}$ konvergen: sehingga $\bigl(f^{-1}\bigr)'(b) =
1/f'(a)$.
\end{proof}

\section{Latihan}

\begin{exercise}[$\star$]\label{exo:b3:residues:1}
Golongkan kesingularannya di $0$ lalu hitunglah \omterm{def:b3:residues:singularities}{residunya}:
\[
\frac{\sin z}{z},\qquad
\frac{\eu^z - 1}{z^2},\qquad
\frac{1}{z(z-1)^2},\qquad
\frac{\cos z}{z^3},\qquad
\eu^{1/z},\qquad
\frac1{\sin z} .
\]
Berikan pula \omterm{def:b3:residues:singularities}{residu} yang ketiga di $z = 1$ dan yang terakhir
di $z = \pi$.
\end{exercise}

\begin{exercise}[$\star$]\label{exo:b3:residues:2}
Untuk $a > 1$ hitunglah, lewat $z = \eu^{\iu t}$:
\[
\int_0^{2\pi}\frac{\dd t}{a + \cos t}
= \frac{2\pi}{\sqrt{a^2 - 1}} .
\]
Periksalah perilaku limitnya saat $a \to 1^+$ dan $a \to \infty$.
\end{exercise}

\begin{exercise}[$\star\star$]\label{exo:b3:residues:3}
Hitunglah dengan kontur setengah lingkaran, sambil membenarkan taksiran busurnya:
\[
\int_\R\frac{x^2}{1 + x^6}\,\dd x = \frac\pi3,
\qquad
\int_\R\frac{\dd x}{(1 + x^2)^{2}} = \frac\pi2
\quad\text{(yang menurunkan kembali \cref{exo:b3:fouriertransform:3})}.
\]
\end{exercise}

\begin{exercise}[$\star\star$]\label{exo:b3:residues:4}
Buktikan, untuk $t \geq 0$ dan $a > 0$:
\[
\int_\R\frac{\cos(tx)}{x^2 + a^2}\,\dd x =
\frac{\pi}{a}\,\eu^{-at},
\]
lalu turunkan transformasi Fourier $x \mapsto \eu^{-a\abs
x}$ lewat pembalikannya --- sambil membandingkannya dengan
\cref{exo:b3:fouriertransform:1}.
\end{exercise}

\begin{exercise}[$\star\star$]\label{exo:b3:residues:5}
Uraikan $f(z) = \dfrac1{(z-1)(z-2)}$ menjadi \omterm{thm:b3:residues:laurent}{deret Laurent} pada masing-masing
ketiga daerahnya $\abs z < 1$, $1 < \abs z < 2$, dan $\abs z >
2$. Mengapa ketiga uraiannya berbeda? Lalu terangkanlah mengapa
koefisien $z^{-1}$ pada uraian kedua dan ketiganya \emph{bukan}
\omterm{def:b3:residues:singularities}{residu} $f$ di $0$ (sebab $f$ tak berkesingularan
di sana), lalu hitunglah \omterm{def:b3:residues:singularities}{residu} $f$ yang sebenarnya, di $1$ dan di
$2$.
\end{exercise}

\begin{exercise}[$\star\star$]\label{exo:b3:residues:6}
(a) Tunjukkan bahwa $\eu^{1/z}$ berkesingularan hakiki di $0$
lalu periksalah Casorati--Weierstrass dengan tangan: selesaikan $\eu^{1/z} =
w$ secara eksplisit untuk sembarang $w \neq 0$, sambil menampilkan solusi
yang sedekat-dekatnya dengan $0$.
(b) Tunjukkan bahwa $\abs{\eu^{1/z}}$ tak terbatas pada setiap
\omterm{def:b3:topology:topology}{persekitaran} tertusuk $0$ padahal $\eu^{1/z}$ tak berkutub:
limit yang manakah yang gagal?
\end{exercise}

\begin{exercise}[$\star\star$]\label{exo:b3:residues:7}
Cacahlah dengan Rouché:
(a) akar $z^7 - 4z^3 + z - 1$ di $\abs z < 1$;
(b) akar $z^4 + 5z + 1$ di $\abs z < 1$ dan di $1 <
\abs z < 2$;
(c) lalu buktikan kembali d'Alembert--Gauss: bahwa polinomial monik berderajat $n$
berakar $n$ kali di suatu cakram yang besar \emph{(bandingkan dengan
$z^n$)}.
\end{exercise}

\begin{exercise}[$\star\star\star$]\label{exo:b3:residues:8}
(Hurwitz) Misalkan $f_n \to f$ secara seragam pada \omterm{def:b3:topology:compact}{kompak}, dengan $f_n \in
\mathcal H(\Omega)$, $\Omega$ \omterm{def:b3:topology:connected}{terhubung}, dan $f \not\equiv 0$.
(a) Tunjukkan bahwa jika semua $f_n$ bebas akar, maka demikian pula $f$.
\emph{(Sebab jika $f(a) = 0$: pakailah asas argumen pada sebuah lingkaran kecil
di sekitar $a$, lalu \cref{thm:b3:holomorphic:weierstrassconv} untuk
mengambil limitnya pada $\int f_n'/f_n$.)}
(b) Tunjukkan bahwa jika semua $f_n$ injektif, maka $f$ injektif atau
konstan. \emph{(Terapkan (a) pada $z \mapsto f_n(z) - f_n(w)$ di
$\Omega\setminus\{w\}$.)}
\end{exercise}

\begin{exercise}[$\star\star\star$]\label{exo:b3:residues:9}
Untuk $n \geq 2$, integralkan $\frac1{1 + z^n}$ atas batas
juringnya $\{0 \leq \arg z \leq \frac{2\pi}n,\ \abs z
\leq R\}$ lalu turunkan
\[
\int_0^{+\infty}\frac{\dd x}{1 + x^n} =
\frac{\pi}{n\,\sin(\pi/n)} .
\]
Periksalah $n = 2$ terhadap $\arctan$, beserta limitnya saat $n \to
\infty$.
\end{exercise}

\begin{exercise}[$\star\star$]\label{exo:b3:residues:10}
Misalkan $f$ sebuah fungsi rasional dengan $\deg(\text{penyebut})
\geq \deg(\text{pembilang}) + 2$. Tunjukkan bahwa jumlah
\emph{semua} \omterm{def:b3:residues:singularities}{residu} $f$ bernilai nol \emph{(integralkan atas
lingkaran yang makin besar)}. Lalu pakailah ini untuk menghitung ulang penguraian
pecahan parsial $\frac1{z(z-1)(z-2)}$ tanpa aljabar
linear.
\end{exercise}

\begin{exercise}[$\star\star\star$]\label{exo:b3:residues:11}
(Lubang kunci: integral pencerminan Euler) Untuk $0 < a < 1$,
hitunglah
\[
I(a) = \int_0^{\infty}\frac{x^{a-1}}{1 + x}\,\dd x =
\frac{\pi}{\sin(\pi a)}
\]
dengan mengintegralkan $f(z) = \frac{z^{a-1}}{1+z} =
\frac{\eu^{(a-1)\log z}}{1 + z}$ (dengan logaritmanya teriris sepanjang
$\R_+$ dan $\arg z \in \intoo0{2\pi}$) atas kontur lubang
kuncinya: keluar sepanjang sisi atas irisannya dari $\varepsilon$ ke
$R$, mengelilingi $C_R$, kembali di bawah irisannya, lalu mengelilingi
$C_\varepsilon$. Benarkanlah: bahwa kedua penggal lurusnya berselisih
sebesar faktor $\eu^{2\iu\pi(a-1)}$, sumbangan lingkarannya
lenyap (sebab $R^{a-1}\cdot R \to 0$ dan
$\varepsilon^{a-1}\cdot\varepsilon \to 0$), dan satu-satunya
kutubnya $z = -1$ berresidu $\eu^{\iu\pi(a - 1)}$. Lalu turunkan pula
$\Gamma(a)\Gamma(1 - a) = \frac\pi{\sin\pi a}$ \emph{(tulislah
$\Gamma(a)\Gamma(1-a) = B(a, 1-a)$ menurut
\cref{pb:b3:lebesgue:1} lalu substitusikan $t =
\frac{x}{1+x}$)}.
\end{exercise}

\begin{exercise}[$\star\star$]\label{exo:b3:residues:12}
(Mencacah akar lewat asas argumen, secara numerik)
Misalkan $P(z) = z^4 + 8z + 1$.
(a) Berapa akarnya di cakram satuannya? \emph{(Pakai Rouché terhadap
$8z + 1$.)}
(b) Berapa akarnya di anulus $1 < \abs z < 3$?
\emph{(Pakai Rouché terhadap $z^4$ pada $\abs z = 3$.)} Lalu tajamkanlah:
tunjukkan bahwa setiap akarnya bermodulus $< 2.1$.
(c) Berapa akarnya di setengah bidang kanannya? \emph{(Cacahlah pada
$\abs z = 2$ lebih dahulu; lalu lacaklah peta sumbu
imajinernya: sebab $P(\iu t) = t^4 + 1 + 8\iu t$ berbagian real positif
di sepanjangnya, sehingga tak ada akar pada sumbunya, dan
ragam argumennya di sepanjangnya dapat dihitung --- lalu simpulkan dengan sebuah
setengah cakram yang besar.)}
\end{exercise}

\section{Soal: \texorpdfstring{$\zeta(2k)$}{zeta(2k)} lewat
kotangen}

\begin{problem}[{Soal akhir pekan --- menjumlahkan $\sum n^{-2k}$
dengan \omterm{def:b3:residues:singularities}{residu}}]\label{pb:b3:residues:1}
Teorema \omterm{def:b3:residues:singularities}{residu} menjumlahkan deret: sebab memasangkan sebuah fungsi rasional
dengan $\pi\cot(\pi z)$, yang kutubnya duduk di bilangan bulat, mengubah
$\sum_{n}f(n)$ menjadi pencacahan \omterm{def:b3:residues:singularities}{residu}. Kita membuktikan metodenya lalu
menghitung $\zeta(2) = \frac{\pi^2}{6}$ dan $\zeta(4) =
\frac{\pi^4}{90}$ --- yakni nilai yang ditemukan lewat deret Fourier pada
Tahun ke-2 dan lewat trace operator pada \cref{ch:b3:spectral}, kini lewat
pengintegralan kontur.

\medskip
\noindent\textbf{Bagian I --- Kernel kotangennya.}
\begin{enumerate}
  \item Tunjukkan bahwa $\pi\cot(\pi z)$ bersifat \omterm{def:b3:residues:singularities}{meromorfik} pada $\C$ dengan
        kutub sederhana tepat di $z = n \in \Z$, yang masing-masing
        berresidu $1$ \emph{(hitunglah
        $\lim_{z\to n}(z-n)\pi\cot\pi z$)}.
  \item Hitunglah awal uraian Laurentnya di $0$:
        \[
        \pi\cot(\pi z) = \frac1z - \frac{\pi^2}{3}\,z -
        \frac{\pi^4}{45}\,z^3 + O(z^5) ,
        \]
        dengan membagi deret pangkat $\cos$ oleh deret $\sin$ (dan benarkanlah
        pembagiannya: sebab $\frac{\sin\pi z}{\pi z}$
        \omterm{def:b3:holomorphic:holo}{holomorfik} dan tak nol di dekat $0$, sehingga
        kebalikannya \omterm{def:b3:holomorphic:holo}{holomorfik}; lalu kenali koefisiennya
        sampai orde $3$).
  \item Misalkan $C_N$ batas bujur sangkar bertitik sudut
        $(\pm1\pm\iu)(N + \frac12)$. Tunjukkan bahwa
        $\abs{\cot(\pi z)} \leq 2$ pada $C_N$ untuk setiap $N \geq
        1$. \emph{(Pada sisi tegaknya, $\cot(\pi(\pm(N +
        \frac12) + \iu y)) = \mp\tan(\iu\pi y)$, yang bermodulus
        $\abs{\tanh(\pi y)} \leq 1$; sedangkan pada sisi mendatarnya
        dengan $\abs y = N + \frac12$, batasilah
        $\abs{\cot(\pi(x\pm\iu y))} \leq \coth(\pi y) \leq
        \coth(\pi/2) < 1.1$.)}
\end{enumerate}

\medskip
\noindent\textbf{Bagian II --- Teorema penjumlahannya.}
\begin{enumerate}[resume]
  \item Misalkan $f$ rasional, \omterm{def:b3:holomorphic:holo}{holomorfik} di bilangan bulatnya,
        dengan $\deg(\text{penyebut}) \geq \deg(\text{pembilang}) + 2$.
        Dengan memakai teorema \omterm{def:b3:residues:singularities}{residu} pada $C_N$ dan batas
        pertanyaan 3, buktikan:
        \[
        \lim_{N\to\infty}\ \sum_{n = -N}^{N} f(n)
        = -\sum_{p\ \text{kutub}\ f}
        \operatorname{Res}\bigl(\pi\cot(\pi z)f(z),\,p\bigr).
        \]
  \item Di manakah argumennya memerlukan syarat derajatnya?
        Tunjukkan lewat contoh (ambillah $f(z) = 1/(z + \frac12)$)
        bahwa untuk peluruhan yang lebih lambat limit setangkupnya
        masih dapat ada sedangkan deret dua sisinya divergen
        --- dan bahwa rumusnya lalu menghitung
        \emph{nilai pokoknya}.
\end{enumerate}

\medskip
\noindent\textbf{Bagian III --- Nilainya.}
\begin{enumerate}[resume]
  \item Terapkan metodenya pada $f(z) = 1/z^2$: di sini $f$ berkutub
        \emph{di} sebuah bilangan bulat, sehingga jalankanlah argumennya
        secara langsung --- integralkan $g(z) = \frac{\pi\cot(\pi
        z)}{z^2}$ atas $C_N$, tunjukkan bahwa integralnya $\to 0$, lalu
        hitunglah $\operatorname{Res}(g, 0)$ dari pertanyaan 2.
        Lalu simpulkan:
        \[
        2\sum_{n\geq1}\frac1{n^2} = \frac{\pi^2}{3},
        \qquad \zeta(2) = \frac{\pi^2}6 .
        \]
  \item Lakukan hal yang sama dengan $g(z) = \frac{\pi\cot(\pi z)}{z^4}$:
        hitunglah $\operatorname{Res}(g, 0)$ lalu turunkan
        $\zeta(4) = \frac{\pi^4}{90}$.
  \item Terangkan pola umumnya: bahwa untuk setiap $k \geq 1$,
        $\zeta(2k)$ adalah $-\frac12$ kali koefisien
        $z^{2k-1}$ pada uraian Laurent $\pi\cot(\pi
        z)$ di $0$ --- yakni kelipatan rasional $\pi^{2k}$.
        Lalu hitunglah $\zeta(6)$ dengan mendorong pembagian pertanyaan 2
        satu langkah lagi. Apa yang dikatakan metodenya tentang
        $\zeta(3)$ --- dan mengapa ia tak mengatakan apa-apa?
\end{enumerate}

\medskip
\noindent\textbf{Bagian IV --- Uraian pecahan parsial
kotangennya.}
\begin{enumerate}[resume]
  \item Tetapkan $w \in \C\setminus\Z$ lalu terapkan metode
        Bagian II pada $f(z) = \dfrac{1}{(z - w)(z + w)}$ ---
        sambil mencatat bahwa $\pi\cot(\pi z)f(z)$ kini berkutub
        sederhana tambahan di $\pm w$, yang \omterm{def:b3:residues:singularities}{residunya} harus turut
        dicacah. Lalu turunkan uraian pecahan parsialnya
        \[
        \pi\cot(\pi w) = \frac1w +
        \sum_{n\geq1}\frac{2w}{w^2 - n^2},
        \]
        dengan deretnya konvergen normal pada himpunan bagian \omterm{def:b3:topology:compact}{kompak}
        $\C\setminus\Z$.
  \item Pulihkan dari uraian ini, dengan menguraikan setiap sukunya
        dalam pangkat $w$ (benarkanlah penukarannya), koefisien
        Laurent yang \emph{sama} seperti pada pertanyaan 2 ---
        sehingga lingkarannya menutup: yakni rumus Euler
        $\sum\frac1{n^2} = \frac{\pi^2}6$ merupakan koefisien
        $w$ pada kedua wajah kotangennya. Lalu bandingkan dengan
        bukti deret Fouriernya (Tahun ke-2) dan bukti tracenya
        (\cref{pb:b3:spectral:1}): yakni tiga teori, satu
        bilangan.
\end{enumerate}

\medskip
\noindent\textbf{Bagian V --- Hasil kali Euler bagi sinusnya.}
Uraian pertanyaan 9 adalah turunan logaritmik sebuah
hasil kali tak berhingga; dan kini kita membuktikan pemfaktoran Euler 1734
secara jujur.
\begin{enumerate}[resume]
  \item Untuk $N \geq 1$ tetapkan $P_N(z) =
        z\prod_{n=1}^{N}\bigl(1 - \frac{z^2}{n^2}\bigr)$.
        Tunjukkan bahwa $P_N$ konvergen, secara seragam pada setiap cakram
        $\bar D(0, R)$, ke sebuah fungsi utuh $P$ yang akarnya
        tepat bilangan bulat dan semuanya sederhana. \emph{(Untuk $n
        \geq 2R$ tulislah faktornya sebagai $\exp\log(1 -
        z^2/n^2)$ dengan logaritma pokok
        \cref{exo:b3:holomorphic:3}, batasilah $\abs{\log(1+u)}
        \leq 2\abs u$ untuk $\abs u \leq \frac12$ lewat
        deretnya, lalu eksponensialkan jumlah logaritmanya yang konvergen
        normal; sedangkan berhingga banyak faktor sisanya
        merupakan polinomial. Lalu simpulkan dengan
        \cref{thm:b3:holomorphic:weierstrassconv}.)}
  \item Tunjukkan bahwa pada $\C\setminus\Z$,
        \[
        \frac{P'(z)}{P(z)} = \frac1z +
        \sum_{n\geq1}\frac{2z}{z^2 - n^2} = \pi\cot(\pi z)
        \]
        \emph{(turunkanlah hasil kali berhingganya, lalu ambil
        limitnya dengan memakai
        \cref{thm:b3:holomorphic:weierstrassconv} dan
        kebebasan akar $P$ di luar $\Z$, lalu kutiplah pertanyaan
        9)}.
  \item Tunjukkan bahwa $Q = \sin(\pi z)/P(z)$ diperluas menjadi fungsi
        utuh yang bebas akar dengan $Q' = 0$, lalu
        simpulkan \emph{hasil kali Euler}:
        \[
        \sin(\pi z) = \pi z\prod_{n\geq1}
        \Bigl(1 - \frac{z^2}{n^2}\Bigr)
        \qquad (z \in \C) .
        \]
  \item (Wallis, 1655) Nilaikan di $z = \frac12$:
        \[
        \frac\pi2 = \prod_{n\geq1}\frac{4n^2}{4n^2 - 1}
        = \lim_{N\to\infty}
        \frac{2\cdot2\cdot4\cdot4\cdots(2N)(2N)}
        {1\cdot3\cdot3\cdot5\cdots(2N-1)(2N+1)} .
        \]
  \item Untuk $\abs z < 1$, uraikan logaritma
        hasil kalinya menjadi deret rangkap (benarkanlah
        penyusunan ulangnya) lalu pulihkan $\zeta(2) =
        \frac{\pi^2}6$ dengan mencocokkan koefisien $z^3$
        pada $\sin(\pi z) = \pi z - \frac{\pi^3}6z^3 + \cdots$
        --- yakni wajah hasil kali bilangan Euler itu.
\end{enumerate}

\medskip
\noindent\textbf{Bagian VI --- Kernel saudaranya.} Kotangennya
bersaudara; dan masing-masing menghargai keluarga deretnya sendiri.
\begin{enumerate}[resume]
  \item Turunkan uraian pertanyaan 9 suku demi suku
        (yang dibenarkan
        \cref{thm:b3:holomorphic:weierstrassconv}) untuk
        memperoleh, secara normal pada \omterm{def:b3:topology:compact}{kompak} $\C\setminus\Z$,
        \[
        \frac{\pi^2}{\sin^2(\pi z)} =
        \sum_{n\in\Z}\frac1{(z - n)^2} .
        \]
  \item Nilaikan di $z = \frac12$:
        $\sum_{m\geq0}\frac1{(2m+1)^2} = \frac{\pi^2}8$;
        lalu pulihkan $\zeta(2)$ sekali lagi dengan memisahkan
        bilangan bulatnya menurut paritas.
  \item Periksalah identitas penggandaan $\tan\theta =
        \cot\theta - 2\cot(2\theta)$ lalu turunkan
        \[
        \pi\tan(\pi z) =
        \sum_{m\geq0}\frac{8z}{(2m+1)^2 - 4z^2} ,
        \]
        secara normal pada \omterm{def:b3:topology:compact}{kompak} yang menghindari $\frac12 + \Z$.
  \item Uraikan di sekitar $0$ (dengan $\abs z < \frac12$; lewat Fubini
        lagi): dengan $\lambda(s) =
        \sum_{m\geq0}(2m+1)^{-s}$,
        \[
        \pi\tan(\pi z) =
        \sum_{k\geq0}8\cdot4^k\,\lambda(2k+2)\,z^{2k+1} ;
        \]
        lalu bandingkan dengan $\tan u = u + \frac{u^3}3 + O(u^5)$ untuk
        memulihkan $\lambda(2) = \frac{\pi^2}8$ dan memperoleh
        $\lambda(4) = \frac{\pi^4}{96}$, lalu periksa silang
        $\zeta(4) = \frac{\pi^4}{90}$ lewat $\lambda(4) = (1 -
        2^{-4})\,\zeta(4)$.
  \item Periksalah $\frac1{\sin\theta} = \cot\frac\theta2 -
        \cot\theta$ lalu turunkan
        \[
        \frac{\pi}{\sin(\pi z)} = \frac1z +
        \sum_{n\geq1}(-1)^n\,\frac{2z}{z^2 - n^2} .
        \]
        Lalu periksalah tandanya terhadap \omterm{def:b3:residues:singularities}{residu}
        $\pi/\sin(\pi z)$ di bilangan bulatnya.
  \item Bacalah koefisien $z$-nya: $\eta(2) =
        \sum_{n\geq1}\frac{(-1)^{n-1}}{n^2} =
        \frac{\pi^2}{12}$, lalu benarkanlah kekonsistenannya
        $\eta(2) = (1 - 2^{1-2})\,\zeta(2)$.
  \item (\omterm{def:b3:topology:interior}{Penutup}) Nilaikan uraian pertanyaan 20 di $z =
        \frac12$ lalu turunkan \emph{rumus Leibniz}
        \[
        \frac\pi4 = 1 - \frac13 + \frac15 - \frac17 +
        \cdots
        \]
        Lalu tutuplah dengan sebuah paragraf pendek: satu kernel per
        aritmetikanya --- kernel mana yang menghargai keluarga deret
        yang mana, dan mengapa semuanya secara struktural buta
        terhadap $\zeta(3)$.

\end{enumerate}
\medskip
\noindent\textbf{Bagian VII --- Daftar harga \omterm{def:b3:complete:complete}{lengkapnya}: bilangan
Bernoulli.}
\begin{enumerate}[resume]
  \item Gabungkan uraian pecahan parsial $\pi
        z\cot(\pi z)$ dengan fungsi pembangkit
        bilangan Bernoulli (yakni $\frac{w}{\eu^w - 1} =
        \sum_n\frac{B_n}{n!}w^n$,
        \cref{pb:b3:holomorphic:1}, Bagian VI): dari
        \[
        \pi z\cot(\pi z) = \iu\pi z +
        \frac{2\iu\pi z}{\eu^{2\iu\pi z} - 1}
        \]
        (buktikanlah identitas ini lebih dahulu), turunkan bentuk tertutupnya
        \[
        \zeta(2k) = (-1)^{k+1}\,
        \frac{(2\pi)^{2k}\,B_{2k}}{2\,(2k)!}
        \qquad (k \geq 1).
        \]
  \item Periksalah rumusnya terhadap $B_2 = \frac16$, $B_4 =
        -\frac1{30}$, dan $B_6 = \frac1{42}$: pulihkan $\zeta(2) =
        \frac{\pi^2}6$ dan $\zeta(4) = \frac{\pi^4}{90}$, lalu
        hitunglah $\zeta(6) = \frac{\pi^6}{945}$.
  \item (Rekursi Euler) Uraikan kedua ruas
        $\bigl(z\cot z\bigr)' = \cot z - z(1 + \cot^2z)$ ---
        atau kuadratkan deret kotangennya secara langsung --- untuk membuktikan
        \[
        \Bigl(k + \frac12\Bigr)\zeta(2k) =
        \sum_{j=1}^{k-1}\zeta(2j)\,\zeta(2k - 2j)
        \qquad (k \geq 2),
        \]
        lalu periksalah bahwa ia menghitung $\zeta(4)$ dari $\zeta(2)$ dan
        $\zeta(6)$ dari $\zeta(2), \zeta(4)$ --- yakni semua nilai zeta
        genap dari satu benih $\frac{\pi^2}6$,
        tanpa pengintegralan yang baru.

\end{enumerate}
\end{problem}
